\chapter{Inflación, estructura y pérdida de las condiciones iniciales}

\section{Omega y las pruebas de la planitud}

La pregunta de un estudiante sobre el gran \(\Omega\) lleva la discusión de nuevo a la ecuación de Friedmann-Robertson-Walker. Para la época actual, mucho después de que terminara el dominio de la radiación, la ley de expansión se escribe con un término de energía del vacío y otro de materia; se omite la radiación porque ahora es muy pequeña:
\begin{equation}
H^2=\left(\frac{\dot a}{a}\right)^2=\frac{8\pi G}{3}\left(\rho_0+\rho_m\right)-\frac{k}{a^2}.
\end{equation}
El término del vacío no se diluye. El de la materia sí:
\begin{equation}
\rho_m=\frac{M}{a^3}.
\end{equation}
\(\Omega\) es la fracción de la ecuación de expansión que corresponde a los distintos componentes.

Dividamos por \(H^2\). Entonces, los términos normalizados suman uno. En una notación regularizada,
\begin{equation}
\begin{aligned}
\Omega_0&=\frac{8\pi G\rho_0}{3H^2},\\
\Omega_m&=\frac{8\pi G\rho_m}{3H^2},\\
\Omega_k&=-\frac{k}{a^2H^2},
\end{aligned}
\end{equation}
y, por tanto,
\begin{equation}
\Omega_0+\Omega_m+\Omega_k=1.
\end{equation}
\editorialnote{El audio y la pizarra son poco nítidos en este tramo, y el conferenciante usa varios nombres para la contribución del vacío. Aquí se regulariza la notación del argumento después de dividir por \(H^2\).}
Si el espacio es plano, \(k=0\), de modo que
\begin{equation}
\Omega_0+\Omega_m=1.
\end{equation}
Si la suma de materia y vacío es mayor que uno, la curvatura es positiva. Si es menor que uno, la curvatura es negativa.

La lógica observacional es entonces sencilla. Se mide la materia ordinaria y la materia oscura. Se mide la expansión de Hubble. Se mide geométricamente la curvatura en escalas muy grandes. Obsérvese también que el universo está acelerando. El signo del término del vacío no viene dado por \(H\) en sí, sino por la tendencia de \(\dot a/a\) hacia un valor constante, lo que significa que el factor de escala se aproxima a la rama exponencial,
\begin{equation}
a(t)\propto e^{Ht}.
\end{equation}
Las distintas mediciones se corroboran entre sí e indican que el universo es casi plano, con una desviación de apenas unos pocos puntos porcentuales.

\begin{classroomqa}
\lecturetimestamp{00:10:50}
\audiencequestion{¿La materia incluye la materia oscura?}
\lecturerresponse{Sí. En esta discusión, el término de materia incluye la materia oscura. También se podría añadir la radiación, pero hoy esa fracción es tan pequeña que puede ignorarse.}
\end{classroomqa}

\begin{classroomqa}
\lecturetimestamp{00:11:11}
\audiencequestion{¿Hay alguna medida de la energía oscura aparte de la constante de Hubble?}
\lecturerresponse{Lo importante no es el valor de \(H\) en sí, sino la manera en que la expansión tiende a seguir una ley exponencial. Junto con la medición de la materia y la medición geométrica de la planitud, esto permite inferir el término del vacío.}
\end{classroomqa}

\enlargethispage{1pt}
\section{De la planitud a la estructura}

¿Por qué es tan plano el universo? La propuesta es la inflación. Que sea plano significa que es muy grande, muy liso y muy homogéneo en escalas de miles de millones, en realidad, decenas de miles de millones, de años luz. La expansión exponencial puede estirar una región hasta que tenga ese aspecto.

Pero una virtud puede llevarse demasiado lejos. Un universo perfectamente liso sigue siendo perfectamente liso. La gravedad no genera grumos a partir de una simetría exacta. Si hoy existe estructura, esta tuvo que originarse en una pequeña inhomogeneidad anterior. Al extrapolar hacia atrás la evolución del universo actual, la inhomogeneidad en la época en que el universo se volvió transparente resulta ser de aproximadamente una parte en \(10^5\):
\begin{equation}
\frac{\delta\rho}{\rho}\sim 10^{-5}.
\end{equation}
El universo era extremadamente liso, pero no perfectamente.

A partir de ahí, la gravedad hace lo que las fuerzas ordinarias no hacen en una habitación llena de aire. Una pequeña sobredensidad atrae materia hacia sí. La región vecina pierde materia, de modo que la subdensidad se vuelve aún menos densa y la sobredensidad, aún más densa. La gravedad refuerza el contraste en lugar de borrarlo.

La pregunta natural es por qué esto no ocurrió antes. La respuesta es la presión. En épocas anteriores, la presión de radiación era grande. Usando
\begin{equation}
p=w\rho,
\end{equation}
los casos habituales son
\begin{equation}
w_{\rm rad}=\frac13,\qquad
w_m=0,\qquad
w_{\rm vac}=-1.
\end{equation}
La presión positiva homogeneiza. Expulsa las perturbaciones y alisa el medio. La materia, tratada como carente de presión en escalas grandes, no ofrece ese tipo de soporte. Mientras hay presión, la gravedad queda contrarrestada. Cuando la presión se vuelve despreciable, la gravedad se impone.

La materia oscura es crucial para el orden de los acontecimientos. Los primeros grumos en colapsar son halos de materia oscura. Después, la materia bariónica ordinaria cae en esos halos y, como los bariones pueden disipar energía, se concentra más cerca del centro. La materia oscura no disipa energía de la misma manera, así que permanece distribuida por el halo. El comienzo de esa contracción es algo anterior a la transparencia, por un factor de aproximadamente diez en tiempo; por eso el fondo cósmico de microondas no muestra directamente su inicio.

Una breve digresión aclara la historia química posterior. Casi todos los elementos más pesados que el helio son producto de la evolución estelar, no de la cosmología del universo temprano. El universo temprano produce, en esencia, hidrógeno, helio y pequeñas trazas de unos pocos núcleos ligeros.

Lo que se observa en el cielo no es la densidad de masa en sí, sino anisotropías de temperatura en el fondo cósmico de microondas. Las fluctuaciones de densidad contribuyen a ellas. También contribuyen las fluctuaciones de velocidad. Las regiones sobredensas tienden a verse ligeramente más frías porque los fotones pierden energía adicional al salir de pozos gravitatorios más profundos, mientras que a esto se suman los corrimientos Doppler debidos al movimiento del gas. El análisis numérico completo es complicado. La conclusión sólida es más sencilla: cuando el universo se volvió transparente, la densidad de energía variaba en aproximadamente una parte en \(10^5\).

\begin{classroomqa}
\lecturetimestamp{00:20:29}
\audiencequestion{¿Por qué la gravedad no hizo crecer los grumos antes?}
\lecturerresponse{Porque antes el universo estaba dominado por la radiación, y esta tenía suficiente presión positiva para homogeneizar el medio. La gravedad solo empezó a dominar el crecimiento de los grumos cuando la radiación perdió importancia frente a la materia.}
\end{classroomqa}

\begin{classroomqa}
\lecturetimestamp{00:24:35}
\audiencequestion{¿No hace falta un mecanismo de disipación para que ocurra el colapso?}
\lecturerresponse{La materia bariónica ordinaria sí necesita disipar energía para asentarse hacia el centro. Pero la materia oscura puede empezar a formar halos sin ese tipo de disipación por colisiones. También existe un tipo de amortiguamiento cosmológico ligado a la propia expansión, que será importante más adelante.}
\end{classroomqa}

\begin{classroomqa}
\lecturetimestamp{00:30:34}
\audiencequestion{¿También había fluctuaciones de velocidad, de modo que las galaxias pudieran adquirir momento angular más adelante?}
\lecturerresponse{Sí. El patrón primordial incluía fluctuaciones de densidad y de velocidad. Lo que se observa en el cielo es un patrón de temperatura producido tanto por las variaciones de densidad como por el movimiento del gas.}
\end{classroomqa}

\section{El inflatón y la meseta}

El mecanismo que se introduce ahora es un campo escalar \(\phi\). Su densidad de energía tiene tres componentes:
\begin{equation}
\rho_\phi=\frac12\dot\phi^2+\frac12(\nabla\phi)^2+V(\phi).
\end{equation}
Contribuyen tanto la variación temporal como la variación espacial. Un campo escalar también puede tener una verdadera energía potencial \(V(\phi)\). Esta parte no es exótica; el campo de Higgs es el ejemplo conocido. Lo que importa aquí es la forma del potencial.

El potencial debe ser bastante plano en cierto intervalo, con una ligera pendiente descendente; luego, en un valor crítico, el fondo cae. No cae exactamente hasta cero, sino hasta un valor muy pequeño. En unidades de Planck, la energía actual del vacío es del orden de
\begin{equation}
10^{-123}.
\end{equation}
En la meseta era mucho mayor, aunque seguía estando muy por debajo de la escala de Planck.

\begin{figure}[t]
\centering
\includegraphics[width=0.78\textwidth]{lecture_08_figure_03.png}
\caption{Potencial cualitativo del inflatón, con una amplia meseta, una caída pronunciada y un mínimo bajo.}
\lectureframe{00:43:23}
\end{figure}

Este gráfico no representa un lugar del espacio, sino un valor del campo. Decir que el universo empezó en lo alto de la meseta solo significa que el valor del campo, en una región muy extensa, se encontraba allí arriba. No importa que \(\phi\) variara algo de un lugar a otro al principio, porque la inflación estira enormemente esas variaciones. Tras suficiente expansión, cualquier región local parece casi uniforme. \editorialnote{Las etiquetas de los ejes no se distinguen en el fotograma de la pizarra, pero la explicación oral deja claro que el eje vertical representa la energía potencial \(V(\phi)\) y el horizontal, el valor del campo \(\phi\).}

Mientras el campo permanezca en la parte amplia y de pendiente suave del potencial, la energía potencial se comporta como energía del vacío. Si \(\phi\) estuviera completamente congelado, se comportaría exactamente como energía del vacío. Dejando de lado los factores numéricos omitidos, la relación que se enuncia en la conferencia es simplemente
\begin{equation}
H\propto \sqrt{V(\phi)},
\end{equation}
y de ahí se obtiene la ley casi exponencial
\begin{equation}
a(t)\propto e^{Ht}.
\end{equation}
Casi todo lo que hoy podemos observar ocurre después de que el campo supera el borde. En esa caída, la gran energía, similar a la del vacío, de la meseta se convierte en radiación, materia ordinaria, materia oscura y la diminuta energía del vacío residual que se observa hoy.

\begin{classroomqa}
\lecturetimestamp{00:44:07}
\audiencequestion{¿Sobre qué espacio está definido \(\phi\), y se siguen usando las ecuaciones de Einstein?}
\lecturerresponse{Sí. \(\phi\) es un campo definido en el espaciotiempo, y se sigue usando la misma ecuación de Friedmann. En este régimen, la densidad de energía del lado derecho proviene principalmente de \(V(\phi)\), siempre que el campo cambie con suficiente lentitud.}
\end{classroomqa}

\begin{classroomqa}
\lecturetimestamp{00:47:04}
\audiencequestion{¿Cómo puede la energía del inflatón, similar a la del vacío, convertirse en materia y radiación?}
\lecturerresponse{Cuando el campo rueda hasta el fondo y oscila, esas oscilaciones son los cuantos del campo. Si los cuantos son inestables, pueden desintegrarse en fotones, electrones, positrones y otras partículas conocidas.}
\end{classroomqa}

\section{Rodamiento lento y fricción de Hubble}

Una vez que el campo supera el borde, ¿por qué no habría de oscilar sin parar? ¿Por qué no habría de ir y venir, quizá incluso volver a subir? La respuesta que se da aquí es un mecanismo de amortiguamiento cosmológico. No es fricción microscópica genuina, pero en la ecuación de movimiento se comporta como tal.

Ignoraremos los gradientes espaciales y trataremos el campo como homogéneo. Entonces, el término cinético y el término potencial se asemejan a un solo grado de libertad con dinámica \(T-V\), pero, como son densidades de energía, deben multiplicarse por el volumen físico \(a^3\). Así se llega a un lagrangiano efectivo homogéneo de la forma
\begin{equation}
L_{\mathrm{hom}}=a^3(t)\left[\frac12\dot\phi^2-V(\phi)\right].
\end{equation}
\editorialnote{La línea superior de la pizarra está parcialmente recortada. Este lagrangiano homogéneo se reconstruye solo en la medida necesaria para obtener la ecuación de movimiento fiable que aparece debajo.}

La ecuación que se distingue claramente en la pizarra es
\begin{equation}
\frac{d}{dt}\left(a^3\dot\phi\right)=-a^3\frac{\partial V}{\partial\phi}.
\end{equation}
Al desarrollar la derivada temporal, se obtiene
\begin{equation}
a^3\ddot\phi+3a^2\dot a\,\dot\phi=-a^3\frac{\partial V}{\partial\phi},
\end{equation}
y al dividir por \(a^3\) resulta
\begin{equation}
\ddot\phi+3H\dot\phi=-\frac{\partial V}{\partial\phi}.
\end{equation}
En la explicación oral se corrige explícitamente el signo menos: la fuerza es el negativo de la derivada de la energía potencial.

\begin{figure}[t]
\centering
\includegraphics[width=0.78\textwidth]{lecture_08_figure_05.png}
\caption{Ecuación de movimiento homogénea del inflatón, con el término de fricción de Hubble \(3H\dot\phi\).}
\lectureframe{00:57:10}
\end{figure}

El término \(3H\dot\phi\) es lo que se denomina fricción de Hubble. Surge por completo al derivar el factor \(a^3\). Sin embargo, desde el punto de vista dinámico se parece al término de arrastre que actúa sobre una roca que cae por el agua. Hay un término de fuerza, uno de aceleración y otro lineal en la velocidad. El campo alcanza rápidamente una velocidad terminal y después se desplaza lentamente por la meseta. Ese es el régimen de rodamiento lento. Mientras la pendiente sea suave y el campo se mueva despacio, la energía potencial cambia lentamente, por lo que la expansión sigue siendo casi exponencial.

Cerca del mínimo, el campo oscila. Esas oscilaciones son los cuantos del inflatón. Si los cuantos son inestables, se desintegran; así es como la energía de la meseta se convierte en el contenido de materia y radiación caliente del universo posterior.

\begin{classroomqa}
\lecturetimestamp{00:50:48}
\audiencequestion{¿Por qué el campo no sigue oscilando eternamente y qué produce el amortiguamiento?}
\lecturerresponse{La propia expansión genera un término parecido a la fricción. El factor \(a^3\) del lagrangiano homogéneo produce el término \(3H\dot\phi\), que se comporta como una fuerza de arrastre viscosa.}
\end{classroomqa}

\section{Fluctuaciones cuánticas y salida del horizonte}

Ahora surge un problema aparente. Si la inflación vuelve al universo extraordinariamente liso, ¿por qué no borra las semillas que necesitan las galaxias? La respuesta es la mecánica cuántica. Las fluctuaciones que sobreviven no son restos de un alisamiento imperfecto: se generan de nuevo.

Olvidemos por un momento la expansión. Aunque un campo fuera casi perfectamente homogéneo, la mecánica cuántica no permitiría fijar con exactitud tanto \(\phi\) como \(\dot\phi\). El estado fundamental sigue teniendo fluctuaciones. Por la misma razón, un oscilador armónico ordinario tiene energía de punto cero. En circunstancias cotidianas, esas fluctuaciones son difíciles de observar porque tienen longitudes de onda cortas y oscilan rápidamente, por lo que las observaciones promedian sus efectos.

La expansión cambia la situación. En coordenadas comóviles, el modo del campo mantiene aproximadamente la misma longitud de onda coordenada, mientras que la longitud de onda física crece con el factor de escala:
\begin{equation}
\lambda_{\rm phys}=a(t)\,\Delta x.
\end{equation}
Durante la inflación, la escala de Hubble física es aproximadamente constante:
\begin{equation}
d_H\sim H^{-1},
\end{equation}
Así, el horizonte comóvil correspondiente se contrae como
\begin{equation}
d_{H,\mathrm{coord}}\sim \frac{1}{aH}.
\end{equation}

La analogía de la cuerda permite entender intuitivamente la congelación. Una onda oscila porque cada segmento de la cuerda es atraído por los segmentos vecinos. Pero, si la longitud de onda llega a ser tan grande como el horizonte, los vecinos relevantes quedan fuera del contacto causal. Desaparece la fuerza restauradora. Cuando un modo supera la longitud del horizonte, se congela.

Esto ocurre una y otra vez. Primero se congelan las longitudes de onda mayores; después, las menores. Incluso un modo que empieza con una longitud de onda diminuta acaba estirándose hasta cruzar el horizonte. Por tanto, la inflación no se limita a borrar la estructura: sigue produciendo variaciones congeladas del campo.

El proceso tampoco se detiene por el mero hecho de que la inflación dure mucho. Los modos congelados antiguos se estiran hasta escalas enormes, pero siempre hay nuevos modos que se congelan. El patrón acumulado alcanza un estado estadísticamente estacionario en el sentido pertinente para la cosmología. De ahí surge el espectro aproximadamente invariante de escala: el mismo proceso de cruce del horizonte se repite en escalas logarítmicas sucesivas y produce aproximadamente la misma amplitud en cada banda de longitudes de onda.

Después de que el campo desciende de la meseta, el horizonte vuelve a crecer. Los modos que habían quedado fuera del horizonte inflacionario temporal vuelven a entrar más adelante. En el lenguaje habitual, salen del horizonte durante la inflación y vuelven a entrar después. Para entonces se han estirado hasta escalas macroscópicas, de modo que, cuando empiezan a oscilar de nuevo, lo hacen como estructuras de tamaño enorme.

\begin{classroomqa}
\lecturetimestamp{01:12:02}
\audiencequestion{¿Qué aspecto tiene el horizonte en el diagrama del espaciotiempo?}
\lecturerresponse{Durante la expansión exponencial, la distancia física al horizonte permanece aproximadamente fija en \(c/H\), pero en un diagrama comóvil esa distancia física fija parece cada vez menor porque la escala de coordenadas se está estirando. En coordenadas comóviles, el horizonte se contrae como \(1/(aH)\).}
\end{classroomqa}

\begin{classroomqa}
\lecturetimestamp{01:15:59}
\audiencequestion{¿Qué pasa con la amplitud cuando se estiran las fluctuaciones? ¿Hay un estado estacionario?}
\lecturerresponse{Los modos individuales se diluyen al estirarse, pero nuevos modos siguen congelándose. El resultado es un patrón estadístico estacionario, aproximadamente invariante de escala, y no la desaparición completa de las fluctuaciones.}
\end{classroomqa}

\section{De las variaciones congeladas del campo a las variaciones de densidad}

Una vez que el campo ha adquirido pequeñas variaciones de un lugar a otro, estas se convierten en diferencias de densidad cuando distintas regiones superan el borde en momentos ligeramente diferentes. Una región un poco más abajo en la pendiente llega primero al precipicio. Si su energía del vacío se convierte en radiación, esta se diluye con la expansión. Una región vecina que aún permanece en la meseta sigue dominada por el vacío y no se diluye. Cuando la segunda región cae, la primera ya ha perdido parte de su energía por corrimiento al rojo.

Por eso, un rodamiento más lento produce contrastes de densidad finales mayores. Si el campo se desplaza muy despacio por la meseta, una pequeña diferencia en su valor corresponde a un mayor desfase temporal entre regiones sucesivas que llegan al precipicio. La región que cae primero tiene más tiempo para diluirse antes de que la región que cae después convierta su energía. Cuanto más suave es la pendiente, menor es la velocidad terminal y mayor la variación final de densidad. Una meseta más alta también tiende a producir fluctuaciones mayores. La anchura de la meseta determina principalmente cuánto dura la inflación.

Los datos de las fluctuaciones restringen con precisión solo una combinación de la altura y la pendiente. Si la escala de energía es suficientemente alta, también debería haber una contribución observable de ondas gravitacionales. Los límites numéricos se dan solo como estimaciones aproximadas. A grandes rasgos, si la densidad de energía en la meseta fuera muy superior a aproximadamente \(10^{-14}\) en unidades de Planck, el fondo cósmico de microondas ya mostraría demasiada estructura debida a ondas gravitacionales; si fuera muy inferior a aproximadamente \(10^{-16}\), habría pocas posibilidades realistas de detectar tal señal. \editorialnote{El conferenciante recalca expresamente que estos números son aproximados y dice que no recuerda sus valores precisos. Aquí se conservan solo como estimaciones de orden de magnitud.}

Las unidades se explican de la manera habitual en unidades naturales. La longitud de Compton de una partícula de masa \(m\) es
\begin{equation}
\lambda_C=\frac{\hbar}{mc}.
\end{equation}
Si \(c=\hbar=1\), la masa y la longitud inversa tienen las mismas unidades. La densidad de energía es energía por unidad de volumen, así que sus unidades son las de la masa elevada a la cuarta potencia:
\begin{equation}
[\rho]=(\mathrm{mass})^4.
\end{equation}

Vista en un solo punto, la acumulación de modos congelados parece un paseo aleatorio. Distintos puntos del espacio siguen paseos aleatorios distintos, y eso es precisamente lo que convierte el ruido cuántico en variación espacial.

\begin{classroomqa}
\lecturetimestamp{01:33:04}
\audiencequestion{¿\(\phi\) sigue un paseo aleatorio, y hace eso irrelevante el punto de partida anterior?}
\lecturerresponse{En cualquier punto fijo, los modos congelados sucesivos se suman con fases aleatorias, así que \(\phi\) sí sigue un paseo aleatorio. Ese mismo resultado también plantea una dificultad: la inflación sigue sustituyendo la información anterior por fluctuaciones posteriores, de modo que resulta difícil recuperar la etapa previa a la inflación.}
\end{classroomqa}

\section{Lo que oculta la inflación}

La teoría tiene éxito en un sentido muy concreto: relaciona las fluctuaciones en un enorme intervalo de escalas mediante un único mecanismo subyacente. Pero ese mismo mecanismo borra lo que ocurrió antes. Las irregularidades anteriores se estiran hasta desaparecer, y las fluctuaciones cuánticas generadas después son las que siguen teniendo relevancia. Por eso la inflación resulta a la vez atractiva y frustrante. Explica la uniformidad y el espectro de fluctuaciones, pero también borra gran parte de la historia anterior.

Esto causa inquietud acerca del punto de partida. Una densidad de energía muy alta no es el problema. El problema es que la energía de la meseta no parece ser ni naturalmente del orden de Planck ni naturalmente diminuta. Es alta en unidades ordinarias, pero sigue siendo pequeña en unidades de Planck. No parece un punto de partida evidente para una teoría fundamental.

En este cuadro sencillo tampoco hay un \(t=0\) literal. Cuando la inflación se prolonga demasiado, se pierde casi toda la información sobre cualquier etapa anterior. Si el número de e-folds apenas superara el mínimo necesario, quizá sobreviviría algún rastro de la etapa previa; pero, en general, la inflación borra el pasado.

La misma cautela se aplica al tratamiento del espacio de de Sitter. La ley exponencial pura es correcta para la rama inflacionaria en tiempos tardíos, pero la solución exacta de de Sitter se describe mejor mediante
\begin{equation}
a(t)\propto \cosh(Ht).
\end{equation}
Esa solución se contrae hasta alcanzar un tamaño mínimo y luego se expande. Queda abierta la cuestión de si debe tomarse literalmente la rama de contracción en cosmología.

\begin{classroomqa}
\lecturetimestamp{01:37:31}
\audiencequestion{¿Podría no haber habido un comienzo y haberse prolongado el comportamiento exponencial indefinidamente hacia el pasado?}
\lecturerresponse{No en este cuadro sencillo. El espacio de de Sitter exacto no es una expansión exponencial eterna que se extienda a todo el tiempo; se parece más a \(\cosh(Ht)\), con una rama de contracción y otra de expansión. Se deja abierta la cuestión de si esa rama anterior debe tomarse literalmente.}
\end{classroomqa}

\section{Ajuste fino, curvatura y límites antrópicos}

En primera aproximación, los parámetros importantes son la altura de la meseta, su pendiente y su anchura \(\Delta\phi\). La pendiente importa porque determina la velocidad terminal y, por tanto, el tamaño de las fluctuaciones. La anchura importa porque ayuda a determinar cuánto dura la inflación. Existe la idea preconcebida de que el desplazamiento de un campo fundamental no debería ser mucho mayor que la masa de Planck, y de que la densidad de energía no debería superar un valor del orden de la unidad en unidades de Planck. Con límites tan aproximados, no todo el espacio de parámetros permite una inflación suficiente. La estimación que se cita aquí es un ajuste fino de unos pocos puntos porcentuales, aproximadamente \(2\%\), para que la pendiente sea lo bastante suave como para obtener entre aproximadamente \(60\) y \(62\) e-folds.

\begin{classroomqa}
\lecturetimestamp{01:46:02}
\audiencequestion{¿Puede entenderse la inflación en términos de un espacio de parámetros, y qué grado de ajuste fino requiere?}
\lecturerresponse{Sí. En primera aproximación, se consideran la altura, la pendiente y la anchura de la meseta. La conferencia da una estimación aproximada: conseguir suficientes e-folds requiere un ajuste fino de unos pocos puntos porcentuales, aproximadamente \(2\%\).}
\end{classroomqa}

Ese ajuste fino moderado se compara luego con el enigma mucho más agudo de la constante cosmológica actual, cuyo orden de magnitud es \(10^{-123}\). La discusión pasa a las cuestiones antrópicas porque ambos problemas pueden replantearse preguntando qué hace falta para que se forme estructura.

En el cuadro newtoniano, una constante cosmológica positiva actúa como una repulsión universal proporcional a la distancia. Si fuera mucho mayor, no tendría que separar hoy las galaxias para resultar fatal. Bastaría con que fuera lo suficientemente grande como para impedir que las fluctuaciones primordiales de \(10^{-5}\) llegaran a condensarse en galaxias. El ejemplo que se da es un valor aproximadamente dos órdenes de magnitud mayor que el observado, que ya produciría ese efecto. Este es el argumento asociado a Weinberg: la constante cosmológica no tiene que ser exactamente cero; basta con que sea lo suficientemente pequeña para permitir la formación de estructura.

La curvatura negativa actúa en la misma dirección. En el lenguaje newtoniano de las conferencias anteriores, \(k<0\) significa que la velocidad de recesión supera la velocidad de escape. Un impulso excesivo hacia afuera también impide que las pequeñas sobredensidades primordiales lleguen a convertirse en galaxias. La inflación diluye esa curvatura. El límite inferior citado es de aproximadamente \(59.5\) e-folds. Luego se toma \(62\) como un ejemplo representativo de una historia inflacionaria que supera holgadamente el umbral.

Esto cambia la pregunta sobre la probabilidad. Si se pregunta qué fracción bruta del espacio de parámetros da lugar a una inflación suficiente, la respuesta quizá sea solo de unos pocos puntos porcentuales. Si se condiciona la pregunta a encontrarse en un universo donde se forman galaxias y existen observadores capaces de plantearla, la probabilidad condicionada puede ser mucho mayor. La discusión termina ahí: no con una solución definitiva, sino con los enigmas del ajuste fino, su reinterpretación antrópica y el reconocimiento de que es aquí donde siguen abiertas las grandes controversias.
